\color{unidarkgreen}{\subsection{Логика очувствления}\label{logic_och}}
\color{black}

\begin{enumerate}

\item
В устройстве предусмотрено два режима очувствления ДЗШ: оперативный и автоматический. Оперативный режим вводится наличием внешнего сигнала \textit{«ЛОЧ/ ООЧ: Оперативный ввод очув.»} и применяется для обеспечения требуемой чувствительности ДЗО при выводе в ремонт мощных присоединений. Автоматический режим вводится наличием внешнего сигнала \textit{«ЛОЧ/ ООЧ: Нормальный режим очув.»} и предназначен для повышения надежности отключения выключателей присоединений при работе ДЗО, в том числе при АПВ шин и при опробовании.

\item
Режим очувствления, вводимый срабатыванием ДЗО, сохраняется на всё время отключения СШ. При наличии сигнала \textit{«ЛО СШ/ ЛО: Фикс. откл.»} режим очувствления удерживается до момента возможного повторного включения СШ при АПВ. Сброс очувствления до повторного включения осуществляется при наличии сигнала \textit{«ЛО СШ/ ЗАПВ: Работа АПВ»}.

\item
Переход в режим очувствления при опробовании осуществляется при отсутствии напряжения на СШ. Очувствление в данном случае сохраняется на всё время опробования.

\item
При АПВ шин режим очувствления вводится на время включения выключателя и последующего отключения в случае неуспешного АПВ. Для исключения излишнего очувствления при постепенном включении всех выключателей предусмотрена возможность вывода режима очувствления при успешном АПВ по истечении выдержки времени, задаваемой уставкой \textbf{«Точув АПВ»}. Вывод режима осуществляется программным ключом \textbf{«Вывод очув при успеш АПВ»}.

\item
Алгоритм работы ЛОЧ приведен на рисунке \ref{pic_log_ochuv:Log_ochuv}. Уставки ЛОЧ приведены в таблице \ref{app_set:loch}.

\medskip
\begin{figure}[H]
\centering
\includegraphics[width=1\textwidth,height=1\textheight,keepaspectratio]{\fbpath/05.01 ДЗО/661.1771 Логика очувствления/img/Log_ochuv.pdf}
\captionof{figure}{Функциональная схема логики очувствления}\label{pic_log_ochuv:Log_ochuv}
\end{figure}
\medskip


\end{enumerate}

